\begin{abstract}
On-device multi-task serving on resource-constrained edge CPUs and microcontrollers is severely bottlenecked by the memory footprint of specialized experts and the latency of switching between them. While Parameter-Efficient Fine-Tuning (PEFT) adapters like LoRA reduce storage, serving dozens of concurrent experts in floating-point precision (FP32/FP16) still occupies substantial on-chip SRAM and incurs severe DRAM-SRAM transfer bandwidth overheads during dynamic weight switching. Standard routing approaches either collapse under heterogeneous streams (heterogeneity collapse) or require sequential micro-batch partitioning that scales latency linearly with the number of tasks. To resolve these deployment-critical bottlenecks, we propose \textbf{Q-SPS} (Quantized Single-Pass Activation-Space Dynamic Blending), a memory-efficient framework that quantizes expert LoRA adapters to low-bitwidth symmetric integers (INT8/INT4) and executes the low-rank additions in pure integer precision within a single, parallel forward pass. To address the fundamental routing-blending execution contradiction of parallel expert serving, we introduce \textbf{CG-Q-SPS} (Conditional Gated Q-SPS), which dynamically bypasses/skips executing expert adapter paths whose routing coefficients fall below a threshold ($\theta=0.01$). This optimization achieves $O(1)$ constant latency for the heavy shared base model backbone while dynamically scaling adapter execution costs to be proportional only to the active experts in the batch. We evaluate our framework using a rigorous, hardware-calibrated analytical simulation of a 12-layer Vision Transformer (ViT-Tiny) across four diverse visual domains (MNIST, Fashion-MNIST, CIFAR-10, SVHN). In our high-fidelity simulations, CG-Q-SPS (INT4) preserves a high simulated joint mean accuracy of \textbf{79.40\%} (recovering \textbf{99.5\%} of the unquantized FP32 ceiling), slashes expert memory footprints by \textbf{87.5\%} (from 2.76 MB down to 0.345 MB), and achieves a projected \textbf{3.97$\times$ physical speedup} (189.1 ms vs. 749.8 ms cumulative latency) over sequential routing state-of-the-art.
\end{abstract}
